\section{Chapter Conclusion}

This chapter covered the theoretical foundations and prior work underlying the proposed system.

The recommendation pipeline builds on two well-established paradigms.
Sequential modelling with DIF-SASRec~\cite{xie2022difsasrec} extends the SASRec Transformer by decoupling content and category attention streams, directly addressing the known limitation of early-fusion side information approaches.
Graph-based embedding with Cleora~\cite{rychalska2021cleora} provides a scalable alternative to explicit graph neural network training, producing co-purchase community representations for 375{,}280 items at low computational cost.

The search pipeline couples dense retrieval via BGE-M3~\cite{chen2024bge} and HNSW~\cite{malkov2018hnsw} approximate search with sparse BM25 retrieval via Tantivy, fused through RRF~\cite{cormack2009rrf}.
Multilingual capability is provided by the NLLB-200~\cite{nllb2022} translation model, which enables Vietnamese queries to be processed without additional encoder fine-tuning.
Visual search over cover images is supported by a CLIP~\cite{radford2021clip} HNSW index.

The survey of existing solutions traces the line from GRU4Rec through SASRec and BERT4Rec to DIF-SASRec, grounding the decision to replace the prior GRU-SeqDQN system with the proposed dual-mode architecture.
The technologies described in Section~\ref{sec:technology} form the implementation substrate on which the proposed solutions in Chapter~4 are built.

